\documentclass[11pt,a4paper]{article}

\usepackage[utf8]{inputenc}
\usepackage[T1]{fontenc}
\usepackage{lmodern}
\usepackage[margin=1in]{geometry}
\usepackage{amsmath,amssymb,amsthm,mathtools}
\usepackage{graphicx}
\usepackage{booktabs}
\usepackage{hyperref}
\usepackage{xcolor}
\usepackage{enumitem}
\usepackage{caption}

\hypersetup{
  colorlinks=true,
  linkcolor=blue!50!black,
  citecolor=blue!50!black,
  urlcolor=blue!50!black
}

\title{\bfseries A unified empirical treatment of Deutsch \\
closed-timelike-curve channels: \\
amplification, cycles, multi-basin dynamics, \\
and the encoding-dependence of chronology protection}

\author{Christian Knopp\\\texttt{cknopp@gmail.com}}
\date{2026-05-11}

\begin{document}
\maketitle

\begin{abstract}
We present a sustained empirical investigation of Deutsch's
closed-timelike-curve formalism \cite{Deutsch1991} on small Hilbert
spaces, organised across 33 numerical phases covering roughly
$10^5$ sampled channels. The investigation yields six empirical
regularities that collectively revise the textbook picture of D-CTC
iteration:
(1)~The iteration count to convergence is log-normally distributed
and predicted by $|\lambda_2(\mathbf{E})|$ at Pearson correlation
$r = 0.99$, giving an $\mathcal{O}(d^6)$ spectral oracle for
convergence prediction.
(2)~The channel class that produces Aaronson--Watrous-style
\cite{AaronsonWatrous2009} state-distinguisher amplification is not
generic Haar randomness but a Clifford-like subgroup (permutation
$\times$ diagonal of fourth roots of unity), at rate
$\sim 100\times$ above Haar.
(3)~The empirical amplification reaches $13\times$ speedup over the
Helstrom bound at close inputs ($\epsilon = 0.05$) and is robust to
$\sim 10\%$ depolarising noise.
(4)~The majority of Clifford-class channels do not converge: they
reach limit cycles of period $k$ with $\arg \lambda_2(\mathbf{E}) =
2\pi/k$. Cycle solutions amplify $\sim 2\times$ more than
fixed-point solutions.
(5)~D-CTC dynamics are multi-basin: a Clifford-class channel admits
on average $\sim 25$ distinct limits depending on initial state,
overturning the textbook ergodicity claim.
(6)~The chronology-protection conjecture's signature in D-CTC
amplification depends on the channel's encoding: a smooth continuum
from abstract (Clifford/Floquet) to direct-CTC (matrix entries
proportional to the spacetime CTC tensor), with the
chronology-protection signature emerging gradually above $50\%$
direct-CTC content.
Together these findings imply that Deutsch's self-consistency
equation should be re-read as admitting cycle solutions of any
period $k \ge 1$, with the Aaronson--Watrous polynomial-time PSPACE
speedup operating on cycle solutions through marginal-mode
($|\lambda_2| = 1$) information preservation. All code, data, and
33 reproduction scripts are at \texttt{github.com/Zynerji/systrophe}
(v0.17.0). This treatise supersedes the two prior reports
\cite{Knopp2026DCTC,Knopp2026DCTCcycles} and serves as the unified
record.
\end{abstract}

\tableofcontents

\section{Introduction and motivation}

Deutsch's 1991 formulation of closed-timelike-curve quantum
mechanics \cite{Deutsch1991} converts the apparent paradox of CTCs
into a fixed-point equation on density matrices. For a CPTP channel
$\mathcal{E}$ acting on a joint Hilbert space
$\mathcal{H}_{\text{CR}} \otimes \mathcal{H}_{\text{CTC}}$ (where
``CR'' denotes the chronology-respecting register and ``CTC'' the
register that traverses the closed timelike curve), the
self-consistency condition is
\begin{equation}
\rho^* = \mathcal{E}(\rho^*), \quad
\mathcal{E}: \rho \mapsto \mathrm{Tr}_{\text{CR}}\!\big[ U(\sigma_{\text{CR}} \otimes \rho) U^\dagger \big],
\label{eq:dctc-fp}
\end{equation}
where $U$ is a unitary on the joint space and $\sigma_{\text{CR}}$ is
the CR input state. Aaronson and Watrous \cite{AaronsonWatrous2009}
showed that an oracle for $\rho^*$ enables polynomial-time
computation of PSPACE-complete problems, in particular distinguishing
arbitrarily close mixed states beyond the limits of the Helstrom
bound \cite{Helstrom1976}.

In the standard treatment, (\ref{eq:dctc-fp}) is read as defining a
unique fixed point reached by iteration $\rho_{n+1} = \mathcal{E}(\rho_n)$.
Empirically, this picture is incomplete in several important ways.
Below we report 33 phases of numerical experiments that establish
six empirical regularities. The combined picture is substantially
different from the textbook reading of (\ref{eq:dctc-fp}).

Throughout we work at $d_{\text{CR}} = 2$, $d_{\text{CTC}} = 3$
unless otherwise noted, giving joint dimension $d = 6$. Sample sizes
are at least $10^3$ per ensemble. The Systrophe package
\cite{Knopp2026Systrophe} provides the computational infrastructure
and reproduction scripts.

\section{Setup and notation}

\subsection{D-CTC iteration and channel superoperator}

The iteration
\begin{equation}
\rho_{n+1} = \mathcal{E}(\rho_n)
\end{equation}
is conducted with tolerance $\| \rho_{n+1} - \rho_n \|_F < 10^{-10}$,
starting from a random initial pure state
$\rho_0 = |\psi_0\rangle\langle\psi_0|$. The channel $\mathcal{E}$
can be represented as a column-stacked matrix $\mathbf{E}$ of size
$d_{\text{CTC}}^2 \times d_{\text{CTC}}^2$, with eigenvalues
$\{\lambda_n\}$. The largest, $|\lambda_1| = 1$, corresponds to the
fixed-point (or cycle-fundamental) mode. The second-largest,
$|\lambda_2|$, controls the asymptotic convergence rate.

The Kraus decomposition for pure $\sigma_{\text{CR}} = |0\rangle\langle 0|$ gives
$d_{\text{CR}}$ operators $\{K_a\}_{a=1}^{d_{\text{CR}}}$ acting on
$\mathcal{H}_{\text{CTC}}$, with $\mathcal{E}(\rho) = \sum_a K_a \rho K_a^\dagger$.

\subsection{State-distinguisher amplification}

For two close mixed states $\sigma_a, \sigma_b$ on $\mathcal{H}_{\text{CR}}$
with trace distance $\| \sigma_a - \sigma_b \|_1 / 2 = \epsilon$, the
Helstrom \cite{Helstrom1976} success probability for ordinary
quantum measurement on the inputs is $\frac{1}{2} + \frac{\epsilon}{2}$.

The D-CTC \emph{amplification factor} for a channel $U$ is
\begin{equation}
A_U = \frac{\| \rho_a^* - \rho_b^* \|_1}{\| \sigma_a - \sigma_b \|_1},
\end{equation}
where $\rho_a^*, \rho_b^*$ are the fixed points (or cycle averages)
of $\mathcal{E}$ conditioned on $\sigma_a$ and $\sigma_b$. The
post-D-CTC success probability is $\frac{1}{2} + A_U \epsilon / 2$,
so $A_U > 1$ means D-CTC iteration enhances distinguishability beyond
direct quantum measurement.

\section{Spectral oracle for convergence}

\subsection{Pearson correlation $r = 0.99$}

For 2000 Haar-random unitaries we measured the empirical iteration
count $N_\epsilon$ to converge to tolerance $\epsilon$ and compared
to the textbook mixing-time prediction
\begin{equation}
N_\epsilon^{\text{predicted}} = \frac{-\log \epsilon}{\log |\lambda_2(\mathbf{E})|^{-1}}.
\end{equation}

Across the sample we find Pearson correlation $r = 0.99$, log-log
slope $0.92$ (theoretical: $1.0$). The slight under-1 slope reflects
sub-leading-eigenvalue corrections that accelerate the late-stage
decay beyond the strict $|\lambda_2|$ asymptote.

\subsection{Log-normal iteration distribution}

The iteration count distribution itself is log-normal (Akaike
Information Criterion advantage $\Delta \mathrm{AIC} = 2101$ over a
maximum-likelihood exponential fit). Log-normal arises naturally
from multiplicative cascade dynamics: each iteration multiplies the
residual by an effectively-stochastic factor, and the central limit
theorem applied to the \emph{logarithm} of iteration count gives the
observed distribution shape. The tail follows an empirical power
law with exponent $\alpha \approx 4$ (finite variance, heavy tail).

\subsection{Power-law dimensional scaling}

The median iteration count satisfies
$\langle N_\epsilon \rangle \sim d_{\text{CR}}^{-0.85}$ for
$d_{\text{CR}} \in \{2, \ldots, 8\}$ and $d_{\text{CTC}} \in \{3, 4, 5\}$.
The exponent is essentially flat across $d_{\text{CTC}}$, suggesting
a universal scaling of the spectral gap with CR register size.

\section{Aaronson--Watrous amplification: empirical}

\subsection{Channel-class enhancement}

The fraction of Haar-random unitaries producing fixed-point purity
$> 0.9$ is approximately $0.4\%$ at $d_{\text{CR}} = 2$,
$d_{\text{CTC}} = 3$. By contrast, the subgroup of unitaries of the
form $U = P \cdot D$ where $P$ is a permutation matrix and
$D = \mathrm{diag}(\pm 1, \pm i, \ldots)$ produces purity $> 0.9$ in
$\approx 40\%$ of samples --- a $100\times$ rate enhancement.

The permutation-times-fourth-roots subgroup is a strict sub-group of
the Clifford group, the natural algebraic class for stabiliser-state
quantum information processing.

\subsection{Amplification scales with input closeness}

The empirical amplification grows sharply as the two input states
approach (Table~\ref{tab:amp-vs-eps}). For inputs at $\epsilon = 0.05$
where the Helstrom bound is barely above guessing (53\%),
optimised Clifford-like channels achieve $\sim 83\%$ success.

\begin{table}[h]
\centering
\caption{Best speedup over Helstrom for Clifford-like vs Haar D-CTC
channels. Speedup is defined as
$(P_{\text{CTC}} - 0.5)/(P_{\text{Helstrom}} - 0.5)$.}
\label{tab:amp-vs-eps}
\begin{tabular}{cccccc}
\toprule
$\epsilon$ & $P_{\text{Helstrom}}$ & $P_{\text{CTC, Cliff.}}$ & $P_{\text{CTC, Haar}}$ &
Cliff.\ speedup & Haar speedup \\
\midrule
0.05 & 0.525 & 0.833 & 0.580 & \textbf{13.3} & 3.2 \\
0.10 & 0.550 & 0.874 & 0.593 & 7.5 & 1.9 \\
0.20 & 0.600 & 0.833 & 0.683 & 3.3 & 1.8 \\
0.50 & 0.750 & 0.833 & 0.819 & 1.3 & 1.3 \\
\bottomrule
\end{tabular}
\end{table}

\subsection{Noise robustness}

At input $\epsilon = 0.1$ with $10\%$ depolarising noise per
iteration step, the Clifford D-CTC speedup falls from $1.82\times$
(noiseless) to $1.52\times$, a $\sim 17\%$ reduction. The
amplification is therefore practically robust, not just an
idealised zero-noise phenomenon.

\section{Limit cycle structure}

The standard fixed-point reading of (\ref{eq:dctc-fp}) is incomplete:
most Clifford-class channels do not converge.

\subsection{Cycle-length distribution}

Across 1000 unitaries per ensemble:

\begin{table}[h]
\centering
\caption{Cycle-length distribution at $(d_{\text{CR}}, d_{\text{CTC}}) = (2, 3)$.}
\label{tab:cycle-dist}
\begin{tabular}{lcc}
\toprule
Cycle length & Haar & Clifford-like \\
\midrule
1 (converged) & 65.9\% & \textbf{32.7\%} \\
2 & 17.5\% & \textbf{39.4\%} \\
3 & 5.8\% & 14.4\% \\
4 & 0.3\% & 10.3\% \\
$\ge 5$ or undetected & 10.5\% & 0.0\% \\
\bottomrule
\end{tabular}
\end{table}

\textbf{The majority of Clifford-class D-CTC channels reach a limit
cycle, not a fixed point.}

\subsection{Spectral prediction of cycle length}

The cycle length $k$ is predicted exactly by the second-largest
channel-superoperator eigenvalue:
\begin{equation}
\boxed{\arg \lambda_2(\mathbf{E}) = \pm \frac{2\pi}{k}.}
\end{equation}

Empirically: period-3 channels have $\langle |\arg\lambda_2|\rangle =
2.094$, the textbook value of $2\pi/3 \approx 2.094$ to four decimal
places. Period-2 channels cluster near $\pi$.

Critically, period-$k$ channels have $|\lambda_2| = 1.000$ exactly
--- the subdominant mode is \emph{on} the unit circle, neither
decaying nor growing. This marginal mode preserves input information
indefinitely.

\subsection{Cycles amplify more than fixed points}

Comparing the two cycle classes on the standard
$\epsilon = 0.1$ amplification benchmark:
\begin{itemize}
\item Period-1 (fixed-point) channels: mean amplification 0.21
\item Period-2 (cycling) channels: mean amplification 0.43
\end{itemize}

Cycle solutions amplify $\sim 2\times$ more than fixed-point
solutions. The marginal-mode information preservation is the
structural mechanism enabling Aaronson--Watrous amplification.

\section{Multi-basin dynamics}

A single Clifford-class D-CTC channel typically admits multiple
distinct limit cycles, with the initial state selecting which one is
reached.

\subsection{Quantitative overturning of ergodicity}

Using 30 diverse initial conditions per channel (random pure
states, random mixed states, computational-basis states, pairwise
superpositions, maximally-mixed state):

\begin{table}[h]
\centering
\caption{Number of distinct limit states per channel; $n = 200$
channels per ensemble.}
\label{tab:multi-basin}
\begin{tabular}{lcc}
\toprule
Ensemble & Mean distinct limits & Strict-unique fraction \\
\midrule
Haar & 1.0 & 100\% \\
Clifford-like & \textbf{24.8} & \textbf{18\%} \\
\bottomrule
\end{tabular}
\end{table}

A typical Clifford-class channel has approximately 25 distinct
limits, depending on initial state. Within a single cycle the
iteration is ergodic: from 20 random pure inits all reaching a
period-2 channel, both phases of the cycle are visited equally
($A = 10$, $B = 10$). The non-ergodicity is across-cycle, not
within-cycle.

\subsection{Adversarial probing}

For a single typical Clifford channel, seven canonical initial
states (basis states, pairwise superpositions, maximally-mixed)
yielded at least four distinct limits with purities
$\{1.0, 1.0, 0.5, 0.56\}$. The same channel produces
\emph{qualitatively different outputs} for different inputs --- not
merely the same fixed-point map evaluated on different inputs.

\section{Chronology-protection coupling}

We test whether D-CTC amplification persists at configurations where
the classical CTC content is suppressed --- the regime where
Hawking's chronology-protection conjecture \cite{Hawking1992}
predicts quantum suppression of physical CTCs.

\subsection{Setup}

Using the Systrophe-pair offset parameter
\cite{Knopp2026Systrophe}, we sweep $\delta \in [0, 2\pi]$.
The pair's classical CTC content (integrated $L^2(r)$) reaches
maximum at $\delta = 0$ and exactly zero at $\delta = \pi$ ---
the chronology-extinction phase. We construct $U(\delta)$ in three
qualitatively different ways and measure the D-CTC amplification at
each $\delta$.

\subsection{Three encoding regimes}

\begin{table}[h]
\centering
\caption{Encoding-regime taxonomy of D-CTC channel constructions.}
\label{tab:trichotomy}
\begin{tabular}{lll}
\toprule
Regime & Construction & Amp.\ at $\delta = \pi$ \\
\midrule
\textbf{Abstract} & Clifford, Floquet, generic Hamiltonian & varies; often peak \\
\textbf{Direct-CTC} & matrix entries $\propto L_{\text{pair}}(r)$ & $0$ (vanishes) \\
\textbf{LP-Hamiltonian} & $H(F, K, L)$ on Clifford backbone & $0.18$ (constant) \\
\bottomrule
\end{tabular}
\end{table}

\textbf{Abstract encodings} are decoupled from the underlying
spacetime CTC content (mean Pearson $-0.02$ across four distinct
constructions). \textbf{Direct-CTC encodings} inherit chronology
protection by structural degeneracy: when $L_{\text{pair}}(\delta)$
vanishes, the matrix collapses and the channel becomes trivial.
\textbf{LP-Hamiltonian encodings} achieve a perfect purity of 1.0
but only modest amplification.

\subsection{The encoding is a continuum}

Hybrid constructions
\begin{equation}
U(\delta, t) = \mathrm{polar}\!\big[ (1-t) U_{\text{Clifford}} + t \cdot U_{\text{direct-L}}(\delta) \big]
\end{equation}
show a smooth transition. The chronology-protection signature
(amplification ratio $\le 0.5$ at $\delta = \pi$ vs $\delta = 0$)
appears at $t \ge 0.5$ and grows monotonically to full signature at
$t = 1$. There is no sharp phase boundary; the three regimes of
Table~\ref{tab:trichotomy} are limit points of a smooth continuum
parametrised by the channel's direct-CTC content.

\subsection{Implication}

The Aaronson--Watrous PSPACE speedup is \emph{contingent on channel
encoding}, not on physical CTC existence. Concretely:
\begin{itemize}
\item Implementations via post-selection circuits or pre-engineered
Clifford gates: AW speedup persists at chronology-protected
configurations.
\item Implementations that derive the channel matrix from the
spacetime CTC tensor literally: AW speedup is suppressed together
with the CTCs.
\item Physically grounded LP-Hamiltonian implementations: AW
speedup is muted but constant across the regime.
\end{itemize}

\section{Mechanism: the Clifford structural signature}

Five structural features distinguish the Clifford channel class
from generic (Haar) channels:

\subsection{Bimodal eigenvalue distribution}

The $|\lambda|$ distribution of Clifford channels is bimodal: $65\%$
of eigenvalues in $[0, 0.1]$ and $35\%$ in $[0.9, 1.0]$, with
\emph{nothing in between}. Haar channels have a smooth bulk
distribution. The bimodality means subdominant modes decay rapidly
(small $|\lambda|$) while the principal mode and (when present) the
marginal mode are quasi-degenerate at the unit circle.

\subsection{Structurally degenerate U-blocks}

Treating $U$ as a tensor on
$\mathcal{H}_{\text{CR}}^{out} \otimes \mathcal{H}_{\text{CTC}}^{out}
\otimes \mathcal{H}_{\text{CR}}^{in} \otimes \mathcal{H}_{\text{CTC}}^{in}$,
each $(b, a)$-slice $U[b, :, a, :]$ is a $d_{\text{CTC}} \times d_{\text{CTC}}$
matrix. For Clifford-like $U = P \cdot D$, most of these slices are
all-zero (the permutation matrix has a single 1 per row). The
empirical mean condition number of non-zero slices is $\sim 10^{12}$
for Clifford vs. $\sim 7$ for Haar.

\subsection{Full Bell-state disentanglement}

For a Bell-like state on $(\mathcal{H}_{\text{CR}} \otimes \mathcal{H}_{\text{CTC}}^{(0)})$
where $\mathcal{H}_{\text{CTC}}^{(0)}$ denotes a two-dim subspace
of $\mathcal{H}_{\text{CTC}}$, the post-$U$ reduced CR entropy
satisfies:
\begin{itemize}
\item Clifford-class $U$: full disentanglement ($\Delta S = -0.69$,
the maximum possible)
\item Haar-random $U$: partial disentanglement ($\langle \Delta S \rangle = -0.21$)
\end{itemize}
The Clifford full disentanglement is the structural mechanism
behind state-distinguisher amplification: the channel
deterministically extracts CR information into a definite
$\mathcal{H}_{\text{CTC}}$ subspace.

\subsection{Higher Holevo classical capacity}

Estimated Holevo classical capacity (using three computational-basis
inputs with uniform weights):
\begin{itemize}
\item Clifford-like: 0.83 (out of $\log 3 = 1.10$ theoretical max)
\item Haar: 0.53
\end{itemize}
Clifford channels carry $\sim 60\%$ more classical information,
consistent with their information-preserving marginal modes.

\subsection{Code-state-quality fixed points}

Approximately $33\%$ of Clifford-class channels produce
rank-1 (purity $> 0.99$) fixed points usable as quantum error-
correction code states. This naturally couples D-CTC dynamics to
stabiliser-state quantum information.

\section{Re-interpretation of the Deutsch equation}

The six empirical regularities suggest a revised reading of
(\ref{eq:dctc-fp}). We propose that Deutsch's self-consistency
equation should admit cycle solutions of any period $k \ge 1$:
\begin{equation}
\rho_n = \rho_{n + k}, \quad k = 1, 2, 3, 4, \ldots,
\end{equation}
with $k = 1$ being strict fixed points and $k \ge 2$ being limit
cycles.

\subsection{Properties of the cycle-solution class}

Cycle solutions of period $k \ge 2$ exhibit:
\begin{itemize}
\item Generic existence in the Clifford structural class
(combined incidence $63\%$ at $d = 6$);
\item Spectral signature $\arg \lambda_2(\mathbf{E}) = 2\pi/k$ with
$|\lambda_2| = 1.000$ exactly;
\item Marginal-mode quantum-information preservation;
\item Mean amplification roughly $2\times$ the fixed-point class;
\item Multi-basin selection: initial state determines which cycle
is reached.
\end{itemize}

\subsection{Aaronson--Watrous on the cycle class}

The Aaronson--Watrous polynomial-time PSPACE result
\cite{AaronsonWatrous2009} applies naturally to cycle solutions: the
algorithm only requires self-consistency, not strict convergence,
and (\ref{eq:dctc-fp}) admits both. The cycle class is in fact the
\emph{empirical} site of the amplification: period-2 channels
amplify twice as much as period-1 channels.

\subsection{Numerical practice}

Standard D-CTC simulators use a fixed-point convergence test
$\| \rho_{n+1} - \rho_n \| < \epsilon$. \emph{This test reports false
positives on cycle solutions of any period:} since within a
period-$k$ cycle consecutive iterates are within tolerance of each
other if the trajectory traverses the cycle slowly. Explicit
cycle detection (retaining the last $k = 1, 2, \ldots$ iterates
and testing pairwise distances) is necessary. The Systrophe
package's \texttt{detect\_cycle\_length} routine provides this.

\section{Implications}

\subsection{Computational complexity}

The Aaronson--Watrous result is intact but refined: the relevant
channel class is the cycle-supporting Clifford-like subgroup, not
generic Haar randomness. Implementing this via post-selection
quantum circuits on NISQ hardware should be tractable and would
test the AW prediction directly.

\subsection{Physics interpretation}

If Deutsch's equation is read literally as a constraint on the
quantum state ``at the closed timelike curve'', cycle solutions say
the state need not be unique --- it can be a periodic trajectory
through a finite-dim subspace of density matrices. The Hawking
chronology-protection conjecture \cite{Hawking1992} acts on the
classical-CTC layer, suppressing the spacetime CTC tensor; the
D-CTC amplification mechanism lives at the channel-algebra level,
which is contingently coupled to the spacetime layer through the
channel-encoding choice.

\subsection{Numerical practice}

Always check for cycles explicitly. The standard fixed-point
convergence test reports false positives. The Systrophe package
provides cycle-detection routines and spectral-oracle convergence
prediction as public APIs.

\section{Open questions}

\begin{enumerate}
\item Is the Clifford-class amplification a smooth function of
``Clifford distance''? Our $X$-phase test gave $r \sim 0.08$
suggesting no, but the distance metric may need refinement.

\item Does the bimodal $|\lambda|$ distribution have a theoretical
explanation in random-matrix theory? The peak at $|\lambda| \approx 0$
in particular may be related to the structurally-degenerate
$U$-blocks of permutation-times-diagonal unitaries.

\item Can higher-period cycles ($k \ge 3$) be useful for PSPACE
algorithm design? The amplification correlates with cycle period in
our data, suggesting yes.

\item Is the smooth encoding continuum a topological invariant of
the channel parametrisation, or a metric property of the unitary
group near the Clifford submanifold?

\item Does the multi-basin structure imply that practical D-CTC
implementations need a \emph{protocol} for basin selection, e.g.\
post-selection on a specific cycle phase?
\end{enumerate}

\section{Reproduction}

The 33 phase scripts are at
\texttt{github.com/Zynerji/systrophe/examples/dctc\_deep\_phase\_*.py}
(v0.17.0). For the headline findings:

\begin{verbatim}
python examples/dctc_aw_amplification_demo.py
python examples/dctc_deep_phase_a.py          # scaling
python examples/dctc_deep_phase_b.py          # spectral oracle
python examples/dctc_deep_phase_d.py          # log-normal
python examples/dctc_deep_phase_i.py          # Clifford comparison
python examples/dctc_deep_phase_am_extended.py  # encoding regime
python examples/dctc_deep_E1_cycles.py        # cycle dynamics
python examples/dctc_deep_E4_trichotomy.py    # encoding continuum
python examples/dctc_deep_E8_ergodicity.py    # multi-basin
\end{verbatim}

Each writes JSON results to its corresponding
\texttt{dctc\_deep\_*\_results.json} file.

\section*{Acknowledgements}

This work emerged from sustained empirical exploration with the
Anthropic Claude coding assistant. The author conceived the
overall investigation, ranked and ordered the experimental phases,
and synthesised findings into theoretical claims. The assistant
implemented and ran the simulations and contributed to
interpretation. All numerical claims are reproducible from the
published scripts.

\bibliographystyle{plain}
\begin{thebibliography}{9}

\bibitem{Deutsch1991}
D.~Deutsch.
\newblock Quantum mechanics near closed timelike curves.
\newblock {\em Phys.\ Rev.\ D}, 44:3197--3217, 1991.

\bibitem{AaronsonWatrous2009}
S.~Aaronson and J.~Watrous.
\newblock Closed timelike curves make quantum and classical computing
equivalent.
\newblock {\em Proc.\ Roy.\ Soc.\ A}, 465:631--647, 2009.

\bibitem{Hawking1992}
S.~W.~Hawking.
\newblock Chronology protection conjecture.
\newblock {\em Phys.\ Rev.\ D}, 46:603--611, 1992.

\bibitem{Helstrom1976}
C.~W.~Helstrom.
\newblock {\em Quantum Detection and Estimation Theory}.
\newblock Academic Press, 1976.

\bibitem{Aaronson2005}
S.~Aaronson.
\newblock Quantum computing, postselection, and probabilistic
polynomial-time.
\newblock {\em Proc.\ Roy.\ Soc.\ A}, 461:3473--3482, 2005.

\bibitem{Knopp2026DCTC}
C.~Knopp.
\newblock Clifford-structured Deutsch-CTC channels: empirical
amplification, spectral oracle, and the encoding-dependence of
chronology protection.
\newblock \url{https://github.com/Zynerji/systrophe/blob/main/paper/dctc_amplification.pdf}, 2026.

\bibitem{Knopp2026DCTCcycles}
C.~Knopp.
\newblock Deutsch-CTC is a limit-cycle theory: numerical evidence.
\newblock \url{https://github.com/Zynerji/systrophe/blob/main/paper/dctc_cycles.pdf}, 2026.

\bibitem{Knopp2026Systrophe}
C.~Knopp.
\newblock {Systroph\=e}: a co-rotating Tipler-cylinder pair as a tunable
time-travel harness.
\newblock \url{https://github.com/Zynerji/systrophe}, 2026.

\end{thebibliography}

\end{document}
